\documentclass[10pt]{article}
\usepackage[ngerman]{babel}
\usepackage[utf8]{inputenc}
\usepackage[T1]{fontenc}
\usepackage{amsmath}
\usepackage{amsfonts}
\usepackage{amssymb}
\usepackage[version=4]{mhchem}
\usepackage{stmaryrd}
\usepackage{bbold}

\title{Hauptprüfung Abiturprüfung 2022 - Leistungsfach }

\author{Alexander Schwarz\\
www.mathe-aufgaben.com}
\date{}


\begin{document}
\maketitle
 Baden-WürttembergJuni 2022

\section*{Aufgabe B2}
Gegeben sind die Geradenschar $g_{a}: \vec{x}=\left(\begin{array}{c}1 \\ 0 \\ -2\end{array}\right)+t \cdot\left(\begin{array}{c}a \\ 3 \\ -1\end{array}\right), t \in \mathbb{R}, a \in \mathbb{R}$ und die Gerade $h: \vec{x}=\left(\begin{array}{c}1 \\ -6 \\ 0\end{array}\right)+s \cdot\left(\begin{array}{c}0 \\ 3 \\ -1\end{array}\right), s \in \mathbb{R}$.\\
a) Beschreiben Sie die besondere Lage der Gerade h im Koordinatensystem. (0,5 VP)\\
Zeigen Sie, dass die Gerade h zur Schar $\mathrm{g}_{\mathrm{a}}$ gehört.\\
Alle Geraden der Schar $g_{a}$ liegen in einer Ebene E .\\
Bestimmen Sie eine Koordinatengleichung der Ebene E.\\
(Teilergebnis: $\mathrm{E}: \mathrm{x}_{2}+3 \mathrm{x}_{3}=-6$ )\\
b) Bestimmen Sie denjenigen Wert von a , für den $\mathrm{g}_{\mathrm{a}}$ die $\mathrm{x}_{2}$-Achse schneidet.\\
(1,5 VP)\\
Es gibt zwei Geraden der Schar $\mathrm{g}_{\mathrm{a}}$, die die Gerade h im Winkel $45^{\circ}$ schneiden. Ermitteln Sie die zugehörigen Werte von a.\\
(2,5 VP)\\
c) Bestimmen Sie eine Gleichung einer Gerade, die von allen Geraden der Schar $g_{a}$ den Abstand $\sqrt{40}$ besitzt und zu allen Geraden der Schar $g_{a}$ windschief verläuft.\\
(2,5 VP)

\section*{Lösungen}
\section*{Aufgabe B2}
a) Beschreiben Sie die besondere Lage der Gerade h im Koordinatensystem.

Da die $\mathrm{x}_{1}$-Koordinate des Richtungsvektors von h gleich Null ist und die $x_{1}$-Koordinate des Stützvektors von h ungleich Null ist, verläuft die Gerade h parallel zur $\mathrm{x}_{2}-\mathrm{x}_{3}$-Ebene.

Zeigen Sie, dass die Gerade h zur Schar $\mathrm{g}_{\mathrm{a}}$ gehört.\\
Es muss einen Wert von geben, so dass der Richtungsvektor von h ein Vielfaches zum Richtungsvektor von $\mathrm{g}_{\mathrm{a}}$ ist.

Dies ist für a = 0 der Fall, hier stimmen die beiden Richtungsvektoren überein.\\
Kontrolle ob der Punkt $\mathrm{P}(1|-6| 0)$ auf der Gerade $g_{0}$ liegt:\\
$\left(\begin{array}{c}1 \\ -6 \\ 0\end{array}\right)=\left(\begin{array}{c}1 \\ 0 \\ -2\end{array}\right)+t \cdot\left(\begin{array}{c}0 \\ 3 \\ -1\end{array}\right)$\\
1.Zeile: 1 = 1\\
2.Zeile: $-6=3 t \Leftrightarrow t=-2$\\
3.Zeile: $0=-2-\mathrm{t} \Leftrightarrow \mathrm{t}=-2$

Da das Gleichungssystem die Lösung $t=-2$ besitzt, liegt $P(1|-6| 0)$ auf $h$.\\
Damit sind die Geraden $h$ und $g_{0}$ identisch.

Bestimmen Sie eine Koordinatengleichung der Ebene E.\\
Um zwei Spannvektoren der Ebene E zu ermitteln, wählt man für a zwei verschiedene Werte, z.B. $\mathrm{a}=0$ und $\mathrm{a}=1$.

Spannvektoren der Ebene: $\overrightarrow{u_{0}}=\left(\begin{array}{c}0 \\ 3 \\ -1\end{array}\right)$ und $\overrightarrow{u_{1}}=\left(\begin{array}{c}1 \\ 3 \\ -1\end{array}\right) \wedge$\\
Normalenvektor von $\mathrm{E}: \overrightarrow{\mathrm{n}}=\left(\begin{array}{c}1 \\ 3 \\ -1\end{array}\right) \times\left(\begin{array}{c}0 \\ 3 \\ -1\end{array}\right)=\left(\begin{array}{l}0 \\ 1 \\ 3\end{array}\right)$

Ansatz Koordinatengleichung: $\mathrm{E}: \mathrm{x}_{2}+3 \mathrm{x}_{3}=\mathrm{d}$\\
Einsetzen des Punkts $\mathrm{Q}(1|0|-2)$ in $\mathrm{E}:-6=\mathrm{d}$\\
Koordinatengleichung: $\mathrm{E}: \mathrm{x}_{2}+3 \mathrm{x}_{3}=-6$\\
b) Bestimmen Sie denjenigen Wert von a, für den $\mathrm{g}_{\mathrm{a}}$ die $\mathrm{x}_{2}$-Achse schneidet.

Gleichung der $x_{2}$-Achse: $k: \vec{x}=\left(\begin{array}{l}0 \\ 0 \\ 0\end{array}\right)+s \cdot\left(\begin{array}{l}0 \\ 1 \\ 0\end{array}\right)$

Gleichsetzen der Geradengleichungen von $k$ und $g_{a}$ :\\
$\left(\begin{array}{l}0 \\ s \\ 0\end{array}\right)=\left(\begin{array}{c}1 \\ 0 \\ -2\end{array}\right)+t \cdot\left(\begin{array}{c}a \\ 3 \\ -1\end{array}\right)$\\
1.Zeile: $0=1+\mathrm{a} \cdot \mathrm{t}$\\
2.Zeile: $\mathrm{s}=3 \mathrm{t}$\\
3.Zeile: $0=-2-\mathrm{t}$

Aus der 3.Zeile: $\mathrm{t}=-2$\\
Aus der 2.Zeile folgt dann $\mathrm{s}=-6$\\
Aus der 1.Zeile folgt dann $a=0,5$.

Ermitteln Sie die zugehörigen Werte von a.\\
$\left.\cos \left(45^{\circ}\right)=\frac{\left|\left(\begin{array}{c}0 \\ 3 \\ -1\end{array}\right) \cdot\left(\begin{array}{c}a \\ 3 \\ -1\end{array}\right)\right|}{\left|\left(\begin{array}{c}0 \\ 3 \\ -1\end{array}\right) \cdot\left(\begin{array}{c}a \\ 3 \\ -1\end{array}\right)\right|} \Leftrightarrow \frac{1}{2} \sqrt{2}=\frac{10}{\sqrt{10} \cdot \sqrt{a^{2}+10}} \right\rvert\, \cdot 2 \cdot \sqrt{10} \cdot \sqrt{a^{2}+10}$\\
$\Leftrightarrow \sqrt{2} \cdot \sqrt{10} \cdot \sqrt{\mathrm{a}^{2}+10}=20 \mid$ quadrieren\\
$\Rightarrow 20 \cdot\left(\mathrm{a}^{2}+10\right)=400$\\
$\Leftrightarrow a^{2}+10=20$\\
$\Leftrightarrow a^{2}=10$\\
$\Rightarrow \mathrm{a}= \pm \sqrt{10}$\\
c) Bestimmen Sie eine Gleichung einer Gerade, die von allen Geraden der Schar $\mathrm{g}_{\mathrm{a}}$ den Abstand $\sqrt{40}$ besitzt und zu allen Geraden der Schar $\mathrm{g}_{\mathrm{a}}$ windschief verläuft.

Alle Geraden der Schar befinden sich in der Ebene E: $x_{2}+3 x_{3}=-6$.\\
Da die gesuchte Gerade von allen Geraden den Abstand $\sqrt{40}$ hat muss die Gerade in einer Ebene F liegen, die parallel zu $E$ ist und von $E$ den Abstand $\sqrt{40}$ hat.

Ansatz für die Koordinatengleichung von $F: x_{2}+3 x_{3}=d$\\
HNF von $F: \frac{x_{2}+3 x_{3}-d}{\sqrt{10}}=0$\\
Für den Punkt $Q(1|0|-2)$ soll gelten: $d(Q, F)=\sqrt{40}$\\
$d(Q, F)=\left|\frac{-6-d}{\sqrt{10}}\right|=\sqrt{40}$\\
Daraus folgt $|-6-d|=20$\\
Eine mögliche Lösung ist $\mathrm{d}=14$.\\
Die gesuchte Gerade liegt in der Ebene F: $x_{2}+3 x_{3}=14$\\
Für den Richtungsvektor $\overrightarrow{\mathrm{u}}=\left(\begin{array}{l}x \\ y \\ z\end{array}\right)$ der gesuchten Geraden gilt:\\
1.) Richtungsvektor ist orthogonal zum Normalenvektor von F:

$$
\left(\begin{array}{l}
0 \\
1 \\
3
\end{array}\right) \cdot\left(\begin{array}{l}
x \\
y \\
z
\end{array}\right)=0 \Leftrightarrow y+3 z=0 \Leftrightarrow y=-3 z
$$

2.) Richtungsvektor ist kein vielfacher Vektor von $\left(\begin{array}{c}a \\ 3 \\ -1\end{array}\right)$ für alle Werte von $a$ :

$$
\left(\begin{array}{c}
a \\
3 \\
-1
\end{array}\right) \neq k \cdot\left(\begin{array}{c}
x \\
-3 z \\
z
\end{array}\right)
$$

Diese Ungleichung wird nur für z = 0 erfüllt.\\
Der Wert von x kann beliebig sein, allerdings nicht $\mathrm{x}=0$.

Möglicher Richtungsvektor der gesuchten Geraden: $\left(\begin{array}{l}1 \\ 0 \\ 0\end{array}\right)$\\
Ein Punkt auf der Ebene F lautet z.B. T(0|14|0), den man als Punkt der gesuchten Geraden wählen kann.

Eine mögliche Gleichung der gesuchten Geraden ist $\vec{x}=\left(\begin{array}{c}0 \\ 14 \\ 0\end{array}\right)+t \cdot\left(\begin{array}{l}1 \\ 0 \\ 0\end{array}\right)$


\end{document}